% GENERATED FILE. Edit canonical lessons, then run npm run generate:books.
% canonical-source-hash: fnv1a64:1febd6c01c37d35a
% canonical-lessons: GU-C357-mulya, GU-C357-sanskar, GU-C357-niti, GU-C357-pratha, GU-C357-apnavvu

\chapter{What We Value}
\label{ch:gu-what-we-value}
\hlchaptermodality{357}

\begin{chapteropening}
\textbf{By the end of this chapter:} \emph{I can say what a family or a village values.}

Complete the last lesson of chapter 357: I can say what a family or a village values.
\end{chapteropening}

\section[\texorpdfstring{\gu{મૂલ્ય} (mūlya)}{મૂલ્ય (mūlya)}]{\gu{મૂલ્ય} (mūlya) — a moral value; worth}

\label{lesson:GU-C357-mulya}



\begin{quote}
Before the new one: say the Gujarati for valid, then the Gujarati for to preserve.
\end{quote}

\subsection*{You'll want to know: \gu{મૂલ્ય}}
\textbf{\gu{મૂલ્ય}} — \emph{mūlya} — \textquotedblleft{}a moral value; worth\textquotedblright{}.

\begin{quote}
\gu{અમારા} \gu{કુટુંબમાં} \gu{સાચું} \gu{બોલવું} \gu{એ} \gu{મોટું} \gu{મૂલ્ય} \gu{છે}.
\end{quote}

\begin{quote}
\emph{amārā kuṭumbmā̃ sācũ bolvũ e moṭũ mūlya che.}
\end{quote}

\begin{quote}
In our family, telling the truth matters a great deal.
\end{quote}

A neuter noun.

\begin{cousinweb}
\textbf{\gu{મૂલ્ય}} (\emph{mūlya}) is a tatsama, borrowed from Sanskrit \emph{mūlya}, \textquotedblleft{}price, worth\textquotedblright{} (Monier-Williams); Gujarati also uses it for what people hold dear.
\end{cousinweb}

\begin{grammarlens}[title={What you've built}]
One more word for giving reasons and opinions.
\end{grammarlens}

\subsection*{Your turn}
\begin{itemize}
\raggedright
  \item \emph{Say it:} \emph{mūlya}
  \item \emph{Say it:} \emph{mūlya}, once more
  \item \emph{Say it:} \emph{amārā kuṭumbmā̃ sācũ bolvũ e moṭũ mūlya che.} — In our family, telling the truth matters a great deal.
  \item \emph{Recall:} say the Gujarati for valid, then the Gujarati for to preserve, then say \emph{mūlya} again
\end{itemize}

\begin{tcolorbox}[breakable,colback=teal!4,colframe=teal!35!black,title={Before you move on}]
What does \gu{મૂલ્ય} mean? (\textquotedblleft{}A moral value; worth\textquotedblright{}.) Say it once more.
\end{tcolorbox}
\section[\texorpdfstring{\gu{સંસ્કાર} (sanskār)}{સંસ્કાર (sanskār)}]{\gu{સંસ્કાર} (sanskār) — upbringing; values instilled in a child}

\label{lesson:GU-C357-sanskar}



\begin{quote}
Before the new one: say the Gujarati for to preserve, then the Gujarati for a moral value.
\end{quote}

\subsection*{You'll want to know: \gu{સંસ્કાર}}
\textbf{\gu{સંસ્કાર}} — \emph{sanskār} — \textquotedblleft{}upbringing; values instilled in a child\textquotedblright{}.

\begin{quote}
\gu{આ} \gu{બાળકમાં} \gu{સારા} \gu{સંસ્કાર} \gu{છે}.
\end{quote}

\begin{quote}
\emph{ā bāḷakmā̃ sārā sanskār che.}
\end{quote}

\begin{quote}
This child has been well brought up.
\end{quote}

A masculine noun.

\begin{cousinweb}
\textbf{\gu{સંસ્કાર}} (\emph{sanskār}) is a tatsama, borrowed from Sanskrit \emph{saṃskāra}, \textquotedblleft{}putting together, forming well, perfecting\textquotedblright{} (Monier-Williams).
\end{cousinweb}

\begin{grammarlens}[title={What you've built}]
One more word for giving reasons and opinions.
\end{grammarlens}

\subsection*{Your turn}
\begin{itemize}
\raggedright
  \item \emph{Say it:} \emph{sanskār}
  \item \emph{Say it:} \emph{sanskār}, once more
  \item \emph{Say it:} \emph{ā bāḷakmā̃ sārā sanskār che.} — This child has been well brought up.
  \item \emph{Recall:} say the Gujarati for to preserve, then the Gujarati for a moral value, then say \emph{sanskār} again
\end{itemize}

\begin{tcolorbox}[breakable,colback=teal!4,colframe=teal!35!black,title={Before you move on}]
What does \gu{સંસ્કાર} mean? (\textquotedblleft{}Upbringing; values instilled in a child\textquotedblright{}.) Say it once more.
\end{tcolorbox}
\section[\texorpdfstring{\gu{નીતિ} (nīti)}{નીતિ (nīti)}]{\gu{નીતિ} (nīti) — a policy}

\label{lesson:GU-C357-niti}



\begin{quote}
Before the new one: say the Gujarati for a moral value, then the Gujarati for upbringing.
\end{quote}

\subsection*{You'll want to know: \gu{નીતિ}}
\textbf{\gu{નીતિ}} — \emph{nīti} — \textquotedblleft{}a policy\textquotedblright{}.

\begin{quote}
\gu{સરકારની} \gu{નવી} \gu{નીતિ} \gu{નાના} \gu{વેપારીઓ} \gu{માટે} \gu{સારી} \gu{છે}.
\end{quote}

\begin{quote}
\emph{sarkārnī navī nīti nānā vepārīo māṭe sārī che.}
\end{quote}

\begin{quote}
The government's new policy is good for small traders.
\end{quote}

A feminine noun.

\begin{cousinweb}
\textbf{\gu{નીતિ}} (\emph{nīti}) is a tatsama, borrowed from Sanskrit \emph{nīti}, \textquotedblleft{}leading, guidance, conduct\textquotedblright{} (Monier-Williams).
\end{cousinweb}

\begin{grammarlens}[title={What you've built}]
One more word for giving reasons and opinions.
\end{grammarlens}

\subsection*{Your turn}
\begin{itemize}
\raggedright
  \item \emph{Say it:} \emph{nīti}
  \item \emph{Say it:} \emph{nīti}, once more
  \item \emph{Say it:} \emph{sarkārnī navī nīti nānā vepārīo māṭe sārī che.} — The government's new policy is good for small traders.
  \item \emph{Recall:} say the Gujarati for a moral value, then the Gujarati for upbringing, then say \emph{nīti} again
\end{itemize}

\begin{tcolorbox}[breakable,colback=teal!4,colframe=teal!35!black,title={Before you move on}]
What does \gu{નીતિ} mean? (\textquotedblleft{}A policy\textquotedblright{}.) Say it once more.
\end{tcolorbox}
\section[\texorpdfstring{\gu{પ્રથા} (prathā)}{પ્રથા (prathā)}]{\gu{પ્રથા} (prathā) — a social practice; a custom}

\label{lesson:GU-C357-pratha}



\begin{quote}
Before the new one: say the Gujarati for upbringing, then the Gujarati for a policy.
\end{quote}

\subsection*{You'll want to know: \gu{પ્રથા}}
\textbf{\gu{પ્રથા}} — \emph{prathā} — \textquotedblleft{}a social practice; a custom\textquotedblright{}.

\begin{quote}
\gu{ગામમાં} \gu{આ} \gu{જૂની} \gu{પ્રથા} \gu{હજી} \gu{ચાલે} \gu{છે}.
\end{quote}

\begin{quote}
\emph{gāmmā̃ ā jūnī prathā hajī cāle che.}
\end{quote}

\begin{quote}
This old practice still goes on in the village.
\end{quote}

A feminine noun.

\begin{cousinweb}
\textbf{\gu{પ્રથા}} (\emph{prathā}) is borrowed from Sanskrit \emph{prathā}, \textquotedblleft{}spreading out, extending\textquotedblright{} (Monier-Williams); Gujarati uses it for a custom that has spread and stayed.
\end{cousinweb}

\begin{grammarlens}[title={What you've built}]
One more word for giving reasons and opinions.
\end{grammarlens}

\subsection*{Your turn}
\begin{itemize}
\raggedright
  \item \emph{Say it:} \emph{prathā}
  \item \emph{Say it:} \emph{prathā}, once more
  \item \emph{Say it:} \emph{gāmmā̃ ā jūnī prathā hajī cāle che.} — This old practice still goes on in the village.
  \item \emph{Recall:} say the Gujarati for upbringing, then the Gujarati for a policy, then say \emph{prathā} again
\end{itemize}

\begin{tcolorbox}[breakable,colback=teal!4,colframe=teal!35!black,title={Before you move on}]
What does \gu{પ્રથા} mean? (\textquotedblleft{}A social practice; a custom\textquotedblright{}.) Say it once more.
\end{tcolorbox}
\section[\texorpdfstring{\gu{અપનાવવું} (apnāvvũ)}{અપનાવવું (apnāvvũ)}]{\gu{અપનાવવું} (apnāvvũ) — to adopt}

\label{lesson:GU-C357-apnavvu}



\begin{quote}
Before the new one: say the Gujarati for a policy, then the Gujarati for a social practice.
\end{quote}

\subsection*{You'll want to know: \gu{અપનાવવું}}
\textbf{\gu{અપનાવવું}} — \emph{apnāvvũ} — \textquotedblleft{}to adopt\textquotedblright{}.

\begin{quote}
\gu{આપણે} \gu{સારી} \gu{ટેવો} \gu{અપનાવવી} \gu{જોઈએ}.
\end{quote}

\begin{quote}
\emph{āpaṇe sārī ṭevo apnāvvī joīe.}
\end{quote}

\begin{quote}
We should take up good habits.
\end{quote}

Past, with \gu{તેણે} (he or she …): \gu{તેણે} \gu{નવી} \gu{રીત} \gu{અપનાવી}, \emph{teṇe navī rīt apnāvī}.

\begin{cousinweb}
\textbf{\gu{અપનાવવું}} (\emph{apnāvvũ}) is probably borrowed from Hindi \emph{apnānā}, \textquotedblleft{}to make one's own\textquotedblright{}, built on Hindi \emph{apnā}, \textquotedblleft{}own\textquotedblright{}, which CDIAL traces to a reconstructed \emph{ātmanaka}, \textquotedblleft{}own\textquotedblright{}, under Sanskrit \emph{ātmán}, \textquotedblleft{}self\textquotedblright{} (CDIAL 1135). Gujarati \textbf{\gu{આપણું}} (\emph{āpṇũ}) comes from the same source. CDIAL lists no Gujarati form of the verb.
\end{cousinweb}

\begin{grammarlens}[title={What you've built}]
One more word for giving reasons and opinions.
\end{grammarlens}

\subsection*{Your turn}
\begin{itemize}
\raggedright
  \item \emph{Say it:} \emph{apnāvvũ}
  \item \emph{Say it:} \emph{apnāvvũ}, once more
  \item \emph{Say it:} \emph{āpaṇe sārī ṭevo apnāvvī joīe.} — We should take up good habits.
  \item \emph{Recall:} say the Gujarati for a policy, then the Gujarati for a social practice, then say \emph{apnāvvũ} again
\end{itemize}

\begin{tcolorbox}[breakable,colback=teal!4,colframe=teal!35!black,title={Before you move on}]
What does \gu{અપનાવવું} mean? (\textquotedblleft{}To adopt\textquotedblright{}.) Say it once more.
\end{tcolorbox}
